\documentclass[11pt]{article}

\usepackage[margin=1in]{geometry}
\usepackage{amsmath,amssymb,mathtools}
\usepackage[T1]{fontenc}
\usepackage{lmodern}
\usepackage{microtype}
\usepackage{fancyhdr}

\pagestyle{fancy}
\fancyhf{}
\cfoot{\thepage}
\setlength{\headheight}{14pt}
\setlength{\parindent}{0pt}
\setlength{\parskip}{6pt}

\newcommand{\qed}{\hfill$\square$}

\begin{document}

\begin{center}
    {\Large\textbf{Proof and the Art of Mathematics}}\\[4pt]
    {\large\textbf{Chapter 4 Exercises}}\\[10pt]
    \textbf{Name:} Hanvesh Aduru\\[4pt]
    \textbf{Due Date:} \underline{\hspace{1.8in}}
\end{center}

\section*{Exercise 4.1}

\textbf{Claim.} For every integer $n\geq 4$, $2^n<n!$.

\textbf{Proof by induction on $n$.}

\textit{Base case:} When $n=4$,
\[
2^4=16<24=4!.
\]
Thus the statement holds for $n=4$.

\textit{Induction hypothesis:} Suppose that for some integer
$k\geq 4$, we have $2^k<k!$. We must prove that
$2^{k+1}<(k+1)!$.

Using the induction hypothesis,
\[
2^{k+1}=2\cdot 2^k<2\cdot k!.
\]
Since $k\geq 4$, we know that $k+1>2$. Therefore,
\[
2\cdot k!<(k+1)k!=(k+1)!.
\]
Combining these inequalities gives
\[
2^{k+1}<(k+1)!.
\]
This completes the induction, so $2^n<n!$ for all
integers $n\geq 4$. \qed

\section*{Exercise 4.3}

Let $(f_n)$ be the Fibonacci sequence, where
$f_0=0$, $f_1=1$, and
$f_{n+2}=f_{n+1}+f_n$. We claim that
\[
f_0+f_1+\cdots+f_n=f_{n+2}-1
\]
for every integer $n\geq 0$.

\textbf{Proof by induction on $n$.}

\textit{Base case:} When $n=0$, the left-hand side is
$f_0=0$. The right-hand side is
\[
f_2-1=(f_1+f_0)-1=1-1=0.
\]
Thus the identity holds for $n=0$.

\textit{Induction hypothesis:} Suppose
\[
f_0+f_1+\cdots+f_k=f_{k+2}-1.
\]
For the next case, add $f_{k+1}$ to both sides:
\begin{align*}
f_0+\cdots+f_k+f_{k+1}
&=(f_{k+2}-1)+f_{k+1}\\
&=f_{k+2}+f_{k+1}-1\\
&=f_{k+3}-1.
\end{align*}
The last equality follows from the Fibonacci recurrence.
This is precisely the desired identity for $n=k+1$.
Therefore, by induction, the formula holds for every
$n\geq 0$. \qed

\section*{Exercise 4.5}

Define an alternative Fibonacci sequence by
\[
a_0=0,\qquad a_1=2,\qquad
a_{n+2}=a_{n+1}+a_n.
\]
The sequence begins
\[
0,2,2,4,6,10,\ldots
\]

The analogue of the Fibonacci sum identity is
\[
\boxed{a_0+a_1+\cdots+a_n=a_{n+2}-2.}
\]

\textbf{Proof by induction on $n$.}

\textit{Base case:} When $n=0$, the left-hand side is
$a_0=0$. The right-hand side is
\[
a_2-2=(a_1+a_0)-2=2-2=0.
\]
Thus the identity holds for $n=0$.

\textit{Induction hypothesis:} Suppose
\[
a_0+a_1+\cdots+a_k=a_{k+2}-2.
\]
For the next case, add $a_{k+1}$ to both sides:
\begin{align*}
a_0+\cdots+a_k+a_{k+1}
&=(a_{k+2}-2)+a_{k+1}\\
&=a_{k+2}+a_{k+1}-2\\
&=a_{k+3}-2.
\end{align*}
The final equality follows from the recurrence relation.
Thus the identity holds for $k+1$, completing the induction.

\textbf{Generalization.}

More generally, suppose a sequence satisfies
\[
a_0=A,\qquad a_1=B,\qquad
a_{n+2}=a_{n+1}+a_n.
\]
Then
\[
\boxed{\sum_{i=0}^{n}a_i=a_{n+2}-B}
\]
for every $n\geq 0$.

To prove this, use induction. For $n=0$,
\[
a_0=A=(A+B)-B=a_2-B.
\]
Now suppose $\sum_{i=0}^{k}a_i=a_{k+2}-B$.
Then
\begin{align*}
\sum_{i=0}^{k+1}a_i
&=a_{k+2}-B+a_{k+1}\\
&=a_{k+3}-B.
\end{align*}
This proves the generalized identity by induction.
The constant subtracted from $a_{n+2}$ is the initial
second term $B$. \qed

\section*{Exercise 4.7}

\textbf{Claim.} Every finite set with $n$ elements has
exactly $2^n$ subsets.

\textbf{Proof by induction on $n$.}

\textit{Base case:} The empty set has exactly one subset:
the empty set itself. Since $2^0=1$, the claim holds
for $n=0$.

\textit{Induction hypothesis:} Suppose every set with
$k$ elements has exactly $2^k$ subsets.

Let $S$ be a set with $k+1$ elements. Choose an element
$x\in S$. Every subset of $S$ either contains $x$ or
does not contain $x$.

The subsets that do not contain $x$ are exactly the
subsets of $S\setminus\{x\}$. Since this set has $k$
elements, the induction hypothesis tells us that there
are $2^k$ such subsets.

Every subset that does contain $x$ can be obtained by
taking a subset of $S\setminus\{x\}$ and adding $x$.
There are also $2^k$ subsets of this type.

These two groups are disjoint and account for every
subset of $S$. Consequently, the total number of subsets
is
\[
2^k+2^k=2\cdot 2^k=2^{k+1}.
\]
Thus the claim holds for $k+1$. By induction, every
finite set with $n$ elements has exactly $2^n$ subsets.
\qed

\section*{Exercise 4.9}

\textbf{Claim.} For every positive integer $n$,
\[
\frac{1}{1\cdot2}+\frac{1}{2\cdot3}
+\cdots+\frac{1}{n(n+1)}
=\frac{n}{n+1}.
\]

\textbf{Proof by induction on $n$.}

\textit{Base case:} When $n=1$,
\[
\frac{1}{1\cdot2}=\frac12=\frac{1}{1+1}.
\]
Therefore, the identity holds for $n=1$.

\textit{Induction hypothesis:} Suppose the identity holds
for $n=k$. Then
\[
\sum_{i=1}^{k}\frac{1}{i(i+1)}
=\frac{k}{k+1}.
\]

For $n=k+1$, we obtain
\begin{align*}
\sum_{i=1}^{k+1}\frac{1}{i(i+1)}
&=\frac{k}{k+1}
  +\frac{1}{(k+1)(k+2)}\\
&=\frac{k(k+2)+1}{(k+1)(k+2)}\\
&=\frac{k^2+2k+1}{(k+1)(k+2)}\\
&=\frac{(k+1)^2}{(k+1)(k+2)}\\
&=\frac{k+1}{k+2}.
\end{align*}
This is exactly the desired formula with $n=k+1$.
Therefore, by induction, the identity holds for every
positive integer $n$. \qed

\section*{Exercise 4.13}

The statement is true in the usual logical sense.
It claims that every elephant in the room that is both
flying and polka-dotted is smoking a cigar.

If there are no flying polka-dotted elephants in the
room, then there is no object that violates the claim.
In logic, a universal statement is considered true when
there are no objects satisfying its conditions. This
is called \textit{vacuous truth}.

The statement does not claim that such elephants
actually exist. If a flying, polka-dotted elephant were
in the room and were not smoking a cigar, that elephant
would be a counterexample, and the statement would be
false.

\section*{Exercise 4.15}

The proposed proof is invalid because it confuses the
finiteness of each partial sum with convergence of the
infinite series.

For every positive integer $n$,
\[
S(n)=\sum_{i=1}^{n}\frac1i
\]
is a finite sum of finite numbers, so $S(n)$ is finite.
However, convergence of an infinite series requires its
partial sums to approach a finite limit. Equivalently,
the partial sums must be bounded above.

The induction step establishes only that
\[
S(n+1)=S(n)+\frac{1}{n+1}
\]
is finite whenever $S(n)$ is finite. It does not
establish a single finite upper bound that works for
every $n$.

In fact, the harmonic series diverges. To see this,
group the terms as follows:
\[
1+\frac12+
\left(\frac13+\frac14\right)+
\left(\frac15+\cdots+\frac18\right)+\cdots
\]
After the first two terms, each block has a number of
terms that is a power of two. In the block
\[
\frac{1}{2^m+1}+\cdots+\frac{1}{2^{m+1}},
\]
there are $2^m$ terms, and each term is at least
$1/2^{m+1}$. Therefore, each such block has sum at
least
\[
2^m\cdot\frac{1}{2^{m+1}}=\frac12.
\]
Since there are infinitely many such blocks, the
partial sums grow without bound. Hence
\[
\sum_{n=1}^{\infty}\frac1n=\infty.
\]
The argument in the proposed proof shows only that
each partial sum is finite, not that the infinite
series converges. \qed

\section*{Exercise 4.17}

We prove that the least-number principle (LNP), common
induction (CI), and strong induction (SI) are equivalent.
We work with the nonnegative integers; the same arguments
apply to any integer starting point after shifting indices.

\subsection*{1. Least-number principle implies common induction}

Assume $P(0)$ is true and that for every $n\geq0$,
$P(n)$ implies $P(n+1)$.

Suppose, for contradiction, that $P(n)$ is false for
at least one nonnegative integer. Let
\[
A=\{n\geq0:P(n)\text{ is false}\}.
\]
By the least-number principle, $A$ has a least element,
say $m$.

Since $P(0)$ is true, $m\neq0$, so $m\geq1$.
By the minimality of $m$, $P(m-1)$ must be true.
The induction step then implies that $P(m)$ is true.
This contradicts the fact that $m\in A$.

Therefore, there are no counterexamples, and $P(n)$
holds for every nonnegative integer $n$. Thus the
least-number principle implies common induction.

\subsection*{2. Common induction implies strong induction}

Suppose $P(0)$ is true and, for every $n\geq0$, if
$P(0),P(1),\ldots,P(n)$ are all true, then $P(n+1)$
is true.

Define a new statement $Q(n)$ as follows:
\[
Q(n):\quad P(j)\text{ is true for every }0\leq j\leq n.
\]

The base case $Q(0)$ is true because it is exactly
the statement $P(0)$.

Now suppose $Q(n)$ is true. By its definition,
$P(0),P(1),\ldots,P(n)$ are all true. The premise of
strong induction therefore implies that $P(n+1)$ is
true. It follows that $Q(n+1)$ is true as well.

By common induction, $Q(n)$ holds for every $n\geq0$.
Since $Q(n)$ includes the assertion that $P(n)$ is true,
we conclude that $P(n)$ holds for every $n\geq0$.
Thus common induction implies strong induction.

\subsection*{3. Strong induction implies the least-number principle}

Let $A$ be a nonempty set of nonnegative integers.
We want to prove that $A$ has a least element.

Suppose, for contradiction, that $A$ has no least
element. We will use strong induction to prove that
no nonnegative integer belongs to $A$.

\textit{Base case:} The number $0$ cannot belong to $A$.
If it did, it would be the least element of $A$.
Therefore, $0\notin A$.

\textit{Induction hypothesis:} Suppose none of the
integers $0,1,\ldots,n$ belongs to $A$.

Consider $n+1$. If $n+1$ belonged to $A$, then none
of the smaller nonnegative integers would belong to
$A$, by the induction hypothesis. Consequently,
$n+1$ would be the least element of $A$, contradicting
our assumption. Hence $n+1\notin A$.

By strong induction, no nonnegative integer belongs
to $A$. This contradicts the assumption that $A$ is
nonempty. Therefore every nonempty set of nonnegative
integers has a least element, proving the least-number
principle.

Since we have shown that the least-number principle
implies common induction, common induction implies
strong induction, and strong induction implies the
least-number principle, all three principles are
equivalent. \qed

\end{document}